\documentclass[12pt]{article}


% Lengths ----------------------------------------------------------------------

% save parindent to a new length, originalparindent
\newlength{\originalparindent}
\setlength{\originalparindent}{\parindent}

% set parskip to bigskipamount for space between paragraphs
\setlength{\parskip}{\bigskipamount}

% set parindent to 0pt for disabling paragraph indentation
\setlength{\parindent}{0pt}


% Packages ---------------------------------------------------------------------

% geometry package for setting page
% size, and for refining page margins
\usepackage[a4paper, hscale=0.85, vscale=0.85]{geometry}

% Set font encoding
\usepackage[T1]{fontenc}

% url package for handlink hyperlinks
\usepackage{url}

% hyperref package for handling in-document links and styling links
\usepackage{hyperref}

% fontspec package to load custom fonts
\usepackage{fontspec}

% xcolor package for 
\usepackage[table]{xcolor}

% secdot package for adding dot after section numbers
\usepackage{secdot}

% ulem package for enabling strikethrough 
\usepackage[normalem]{ulem}

% tocloft package for disabling
% bold font in the table of contents
\usepackage{tocloft}

% titletoc package for adding a dot after
% section numbers in the table of contents
\usepackage[dotinlabels]{titletoc}

% setspace package for altering linespread in tables
\usepackage{setspace}

% float package for placing tables
% and figure at exact position
\usepackage{float}

% caption package
% for caption styling 
\usepackage{caption}

% colortbl package for colored tables
\usepackage{colortbl}

% tikz package for inline code
% styling, and horizontal rules
\usepackage{tikz}

% verbatim package for verbatim
% environment in code block environments
\usepackage{verbatim}

% mdframed environment for custom
% code blocks and custom quotes
% (common options for all mdframed based
% environments are set at package loading)
\usepackage[framemethod=tikz,%
    innerleftmargin=0.5\originalparindent,%
    innerrightmargin=0.5\originalparindent,%
    skipabove=0.4\baselineskip,%
    skipbelow=0.4\baselineskip,%
    innertopmargin=0.4\baselineskip,%
    innerbottommargin=0.4\baselineskip]{mdframed}

% tabu package for
% easier tabular styling
\usepackage{tabu}


% Package Setups ---------------------------------------------------------------

% setup for hyperref package:
%   enabled pdf bookmarks,
%   setting link styles
\hypersetup{bookmarks=true,%
    bookmarksnumbered=true,%
    pdfencoding=unicode,%
    colorlinks=true,%
    pdfborder={0 0 0},%
    linkcolor=black,%
    menucolor=black,%
    citecolor=mdhyperlinkcolor,%
    urlcolor=mdhyperlinkcolor,%
    filecolor=mdhyperlinkcolor}

% setup for tikz package:
%   load library for fancy hrlues
\usetikzlibrary{decorations.pathreplacing}

% setup url package:
%   set url font to sans serif instead of teletype
\urlstyle{sf}

% Font settings ----------------------------------------------------------------

\usepackage{sourcesanspro}
\usepackage{sourcecodepro}

% set document default font to Source Sans Pro and its variants
\setmainfont[Mapping=tex-text,%
    ItalicFont=Source Sans Pro Light Italic,%
    BoldFont=Source Sans Pro,%
    BoldItalicFont=Source Sans Pro Italic]{Source Sans Pro Light}

% set sans serif font to Source Sans Pro and its variants (just in case)
\setsansfont[Mapping=tex-text,%
    ItalicFont=Source Sans Pro Light Italic,%
    BoldFont=Source Sans Pro,%
    BoldItalicFont=Source Sans Pro Italic]{Source Sans Pro Light}

% set monospace font to Source Code Pro and its variants
\setmonofont[Mapping=tex-text,%
    ItalicFont=Source Code Pro ExtraLight,%
    BoldFont=Source Code Pro]{Source Code Pro Light}


% Color definitions ------------------------------------------------------------

\definecolor{mdfancyhlinecolor}{HTML}{CCCCCC}
\definecolor{mdsimplehlinecolor}{HTML}{DDDDDD}
\definecolor{mdhyperlinkcolor}{HTML}{4183C4}
\definecolor{mdinlinecodeboxbackgroundcolor}{HTML}{F8F8F8}
\definecolor{mdinlinecodeboxframecolor}{HTML}{DDDDDD}
\definecolor{mdblockquotelinecolor}{HTML}{DDDDDD}
\definecolor{mdalternatingtablerowcolor}{HTML}{F8F8F8}
\definecolor{mdtableframecolor}{HTML}{DDDDDD}
\definecolor{mdimgboxcolor}{HTML}{DDDDDD}


% Styling table of contents ----------------------------------------------------

% set dot fill style
\renewcommand{\cftsecdotsep}{\cftdotsep}
\renewcommand{\cftsubsecdotsep}{\cftdotsep}
\renewcommand{\cftsubsubsecdotsep}{\cftdotsep}
\renewcommand{\cftparadotsep}{\cftdotsep}
\renewcommand{\cftsubparadotsep}{\cftdotsep}
\renewcommand{\cftsecleader}{\cftdotfill{\cftsecdotsep}}
\renewcommand{\cftsubsecleader}{\cftdotfill{\cftsubsecdotsep}}
\renewcommand{\cftsubsubsecleader}{\cftdotfill{\cftsubsubsecdotsep}}
\renewcommand{\cftparaleader}{\cftdotfill{\cftparadotsep}}
\renewcommand{\cftsubparaleader}{\cftdotfill{\cftsubparadotsep}}

% set section font style
\renewcommand\cftsecfont{\normalfont}
\renewcommand\cftsecpagefont{\normalfont}

% set indentation of toc entries
\newlength{\mycftsecindent}
\newlength{\mycftsubsecindent}
\newlength{\mycftsubsubsecindent}
\newlength{\mycftparaindent}
\newlength{\mycftsubparaindent}
\setlength{\mycftsecindent}{0.5\cftsecindent}
\setlength{\mycftsubsecindent}{0.5\cftsubsecindent}
\setlength{\mycftsubsubsecindent}{0.5\cftsubsubsecindent}
\setlength{\mycftparaindent}{0.5\cftparaindent}
\setlength{\mycftsubparaindent}{0.5\cftsubparaindent}
\setlength{\cftsecindent}{\mycftsecindent}
\setlength{\cftsubsecindent}{\mycftsubsecindent}
\setlength{\cftsubsubsecindent}{\mycftsubsubsecindent}
\setlength{\cftparaindent}{\mycftparaindent}
\setlength{\cftsubparaindent}{\mycftsubparaindent}
\addtolength{\cftsecnumwidth}{0.3em}
\addtolength{\cftsubsecnumwidth}{0.3em}
\addtolength{\cftsubsubsecnumwidth}{0.3em}
\addtolength{\cftparanumwidth}{0.3em}
\addtolength{\cftsubparanumwidth}{0.3em}

% set parskip between toc entries
\newlength{\mycftbeforeskip}
\setlength{\mycftbeforeskip}{0.5\cftbeforesecskip}
\setlength{\cftbeforesecskip}{\mycftbeforeskip}
\setlength{\cftbeforesubsecskip}{\mycftbeforeskip}
\setlength{\cftbeforesubsubsecskip}{\mycftbeforeskip}
\setlength{\cftbeforeparaskip}{\mycftbeforeskip}
\setlength{\cftbeforesubparaskip}{\mycftbeforeskip}

% set table of contents depth to 5
\setcounter{tocdepth}{5}


% Section styling --------------------------------------------------------------

% add dots after section numbers (secdot package)
\sectiondot{section}
\sectiondot{subsection}
\sectiondot{subsubsection}
\sectiondot{paragraph}
\sectiondot{subparagraph}

% changing the style of \paragraph and \subparagraph titles, so
% text after \paragraph and \subparagraph are broken into new lines
\makeatletter
    \renewcommand\paragraph{%
        \@startsection{paragraph}{4}{0mm}%
            {-\baselineskip}%
            {.3\baselineskip}%
            {\normalfont\normalsize\bfseries}}
    \renewcommand\subparagraph{%
        \@startsection{subparagraph}{5}{0mm}%
            {-\baselineskip}%
            {.3\baselineskip}%
            {\normalfont\normalsize\bfseries}}
\makeatother

% set section number up to level 5
\setcounter{secnumdepth}{5}

% add a dot after section
% numbers in the pdf bookmarks
% http://tex.stackexchange.com/questions/150983/add-dot-to-the-end-of-section-numbering-in-pdf-bookmarks
\makeatletter
\renewcommand{\Hy@numberline}[1]{#1. }
\makeatother


% Paragraph styling ------------------------------------------------------------

% prevent widows and orphans
\widowpenalty=10000
\clubpenalty=10000

% prevent overfull lines
\sloppy


% New commands -----------------------------------------------------------------

% mdtitle command for document title
% (not listend in the table of contents)
\newcommand{\mdtitle}[1]{{\LARGE\textbf{#1}}}

% mdtableofcontents command for custom styled table of contents
\newcommand{\mdtableofcontents}{{\setlength{\parskip}{0pt}\tableofcontents}}

% mdsimplehrule command for a simple horizontal rule
\newcommand{\mdsimplehrule}{%
    \nopagebreak\begin{tikzpicture}%
        \path[draw, mdsimplehlinecolor] (0, 0) -- (\textwidth{}, 0);%
    \end{tikzpicture}%
}

% mdfancyhrule command for a fancy horizontal rule
\newcommand{\mdfancyhrule}{%
    \nopagebreak\begin{tikzpicture}%
        \pgfdeclaredecoration{fancyhrule}{initial}{%
            \state{initial}[width=4.25pt]%
            {%
                \fill[mdfancyhlinecolor] (0pt, 0pt) -- (3pt, 3pt) -- (4.25pt, 3pt) -- (1.25pt, 0pt) -- cycle;%
            }%
            \state{final}%
            {%
                \pgfpathmoveto{\pgfpointdecoratedpathlast}%
            }%
        }%
        \path[decorate, decoration=fancyhrule] (0, 0) -- (\textwidth, 0);%
    \end{tikzpicture}%
%    \nopagebreak\begin{tikzpicture}[decoration={border, angle=45, segment length=4pt, amplitude=4pt}, thick]%
%        \path[postaction={decorate, draw}, mdfancyhlinecolor] (0, 0) -- (\textwidth{}, 0);%
%    \end{tikzpicture}%
}

% mdinlinecode command for including code snippets inline
% (fake verbatim, so all special character should be escaped,
% or textmode equivalents of special characters should be used)
\newcommand{\mdinlinecode}[1]{%
    \begin{tikzpicture}[baseline=0ex]%
        \node[anchor=base,%
            text height=1em,%
            text depth=1ex,%
            inner ysep=0pt,%
            draw=mdinlinecodeboxframecolor,%
            fill=mdinlinecodeboxbackgroundcolor,%
            rounded corners=2pt] at (0,0) {\footnotesize\texttt{#1}};%
    \end{tikzpicture}%
}

% bfdescriptionlabel command for changing the description
% label style in the mdbfdescription environment
\newcommand{\bfdescriptionlabel}[1]{%
    \hspace{\labelsep}\normalfont{\textbf{#1:}}%
}

% codedescriptionlabel command for changing the description
% label style in the mdcodedescription environment
\newcommand{\codedescriptionlabel}[1]{%
    \hspace{\labelsep}\normalfont{\mdinlinecode{#1}}%
}

% mdimgbox for a frame around figures,
% but it can be used for anything
\newcommand{\mdimgbox}[1]{%
    \setlength{\fboxsep}{0pt}%
    \setlength{\fboxrule}{0.4pt}%
    \fcolorbox{mdimgboxcolor}{white}{#1}%
}

% New environments -------------------------------------------------------------

% mdcodeblock environment for including code blocks
% (based on mdframed, breaks between pages)
\newmdenv[font=\footnotesize,%
linewidth=0.4pt,%
roundcorner=2pt,%
linecolor=mdinlinecodeboxframecolor,%
backgroundcolor=mdinlinecodeboxbackgroundcolor,%
settings={\setlength{\parindent}{0pt}}]{mdcdblk}
\newenvironment{mdcodeblock}{\endgraf\verbatim}{\endverbatim}
\BeforeBeginEnvironment{mdcodeblock}{\begin{mdcdblk}}
\AfterEndEnvironment{mdcodeblock}{\end{mdcdblk}}

% mdnonbreakcodeblock environment for including code blocks
% (based on mdframed, doesn't break between pages)
\newmdenv[font=\footnotesize,%
linewidth=0.4pt,%
roundcorner=2pt,%
linecolor=mdinlinecodeboxframecolor,%
backgroundcolor=mdinlinecodeboxbackgroundcolor,%
nobreak=true,%
settings={\setlength{\parindent}{0pt}}]{mdnonbreakcdblk}
\newenvironment{mdnonbreakcodeblock}{\endgraf\verbatim}{\endverbatim}
\BeforeBeginEnvironment{mdnonbreakcodeblock}{\begin{mdnonbreakcdblk}}
\AfterEndEnvironment{mdnonbreakcodeblock}{\end{mdnonbreakcdblk}}

% mdblockquote environment for custom styled blockquotes
% (based on mdframed, breaks between pages)
\newmdenv[linewidth=3pt,%
linecolor=mdblockquotelinecolor,%
topline=false,%
rightline=false,%
bottomline=false,%
settings={\setlength{\parindent}{0pt}}]{mdblockquote}

% mditemize environment for
% custom styled unordered lists
\newenvironment{mditemize}%
    {\begin{itemize}
        \setlength{\parskip}{0.5\smallskipamount}}%
    {\end{itemize}}

% mdenumerate environment for
% custom styled enumerated lists
\newenvironment{mdenumerate}%
    {\begin{enumerate}
        \setlength{\parskip}{0.5\smallskipamount}}%
    {\end{enumerate}}

% mdbfdescription environment for
% custom (bold) styled description lists
\newenvironment{mdbfdescription}%
{\renewcommand{\descriptionlabel}{\bfdescriptionlabel}%
    \begin{description}%
    \setlength{\itemindent}{\parindent}%
    \setlength{\parskip}{0.5\smallskipamount}}%
{\end{description}}

% mdcodedescription environment for
% custom (mdinlinecode) styled description lists
\newenvironment{mdcodedescription}%
{\renewcommand{\descriptionlabel}{\codedescriptionlabel}
    \begin{description}%
    \setlength{\itemindent}{\parindent}%
    \setlength{\parskip}{0.5\smallskipamount}}%
{\end{description}}


% Document start ---------------------------------------------------------------

\begin{document}

\mdtitle{GitHub Markdown style \LaTeX}

\mdfancyhrule

\mdtableofcontents

\mdfancyhrule

\section{Headings}

Document title with \mdinlinecode{\textbackslash{}mdtitle\{\}}.

Standard \LaTeX section headings can be used (level 1--5).

\subsection{Subsection}

\subsubsection{Subsubsection}

\paragraph{Paragraph}

\subparagraph{Subparagraph}

And these also appear nicely in the table of contents.

\mdsimplehrule

\section{Emphasis}

Emphasis, aka \emph{italics}, with standard \mdinlinecode{\textbackslash{}emph\{\}}.

Strong emphasis, aka \textbf{bold}, with standard \mdinlinecode{\textbackslash{}textbf\{\}}.

\emph{\textbf{Combined emphasis}} with standard \mdinlinecode{\textbackslash{}emph\{\}} and \mdinlinecode{\textbackslash{}textbf\{\}}.

\sout{Strikethrough} with \mdinlinecode{\textbackslash{}sout\{\}} from the \mdinlinecode{ulem} package.

\mdsimplehrule

\section{Lists}

Unordered list with the \mdinlinecode{mditemize} environment:

\begin{mditemize}
\item foo
\item bar
\end{mditemize}

Enumeration, aka ordered list with the \mdinlinecode{mdenumerate} environment:

\begin{mdenumerate}
\item foo
\item bar
\end{mdenumerate}

Bold style description list with the \mdinlinecode{mdbfdescription} environment:

\begin{mdbfdescription}
\item[foo] bar
\item[R. Daneel Olivaw] A fictional robot created by Isaac Asimov.
\end{mdbfdescription}

Inline code description list with the \mdinlinecode{mdcodedescription} environment

\begin{mdcodedescription}
\item[void *] A special type of pointer that can be pointed at objects of any data type.
\item[R. Daneel Olivaw] a fictional robot created by Isaac Asimov.
\end{mdcodedescription}

\emph{Note: the description label inside the \mdinlinecode{mdcodedescription} environment is fake verbatim, so everything must be escaped inside it or textmode equivalents should be used of special characters. (See \hyperref[sec:mdinlinecodedef]{\mdinlinecode{\textbackslash{}mdinlinecode\{\}}'s definition}.)}

\newpage

Combined lists are also available. The style modifications don't affect nesting:

\begin{mdenumerate}
\item This supposed to be an example with custom styled lists.
\begin{mditemize}
\item unordered
\item list
\end{mditemize}
\item Let's see a nested description list.
\begin{mdbfdescription}
\item[The Doctor] So, the year 5 billion. The Sun expands, the Earth gets roasted.
\item[Rose] That was our first date.
\end{mdbfdescription}
\end{mdenumerate}

\mdsimplehrule

\section{Links}

Links with description with \mdinlinecode{\textbackslash{}href\{url\}\{description\}}:

\href{http://en.wikipedia.org/wiki/Three_Laws_of_Robotics}{Three Laws of Robotics - Wikipedia, the free encyclopedia}

URLs with \mdinlinecode{\textbackslash{}url\{\}}:

\url{http://en.wikipedia.org/wiki/Three_Laws_of_Robotics}

\mdsimplehrule

\section{Inline code and code blocks}

Inline code with \mdinlinecode{\textbackslash{}mdinlinecode\{inline code\}}:

Lorem ipsum dolor \mdinlinecode{sit} amet, consectetur adipiscing elit. Phasellus quis auctor \mdinlinecode{nunc}. Nam ac egestas lectus, ac tincidunt purus. Aliquam viverra dictum est, ut pretium magna tristique eget. Donec venenatis \mdinlinecode{purus} \mdinlinecode{eget} justo malesuada suscipit.

\emph{Note: \mdinlinecode{\textbackslash{}mdinlinecode\{\}} is fake verbatim, so everything must be escaped inside it or textmode equivalents should be used of special characters.}

Code blocks with the \mdinlinecode{mdcodeblock} environment (the example shows \mdinlinecode{\textbackslash{}mdinlinecode\{\}}'s definition):

\label{sec:mdinlinecodedef}

\begin{mdcodeblock}
\newcommand{\mdinlinecode}[1]{%
    \begin{tikzpicture}[baseline=0ex]%
        \node[anchor=base,%
            text height=1em,%
            text depth=1ex,%
            inner ysep=0pt,%
            draw=mdinlinecodeboxframecolor,%
            fill=mdinlinecodeboxbackgroundcolor,%
            rounded corners=2pt] at (0,0) {\footnotesize\texttt{#1}};%
    \end{tikzpicture}%
}
\end{mdcodeblock}

Or with the \mdinlinecode{mdnonbreakcodeblock} environment, which won't break between pages.

\mdsimplehrule

\section{Figures}
A framed figure (\mdinlinecode{\textbackslash{}mdimgbox\{\}}) with left aligned caption:

\begin{mdcodeblock}
\begin{figure}[H]
    \captionsetup{justification=justified, singlelinecheck=false}
    \mdimgbox{\textcolor{white!0}{\rule{0.333\textwidth}{0.333\textwidth}}}
    \caption{\url{http://en.wikipedia.org/wiki/Deus_Irae}}
\end{figure}
\end{mdcodeblock}

\begin{figure}[H]
    \captionsetup{justification=justified, singlelinecheck=false}
    \mdimgbox{\textcolor{white!0}{\rule{0.333\textwidth}{0.333\textwidth}}}
    \caption{\url{http://en.wikipedia.org/wiki/Deus_Irae}}
\end{figure}

\mdsimplehrule

\section{Tables}

Table created with the help of \mdinlinecode{colortbl} and the standard table environments:

\begin{mdcodeblock}
{\linespread{1.9}
\begin{table}[H]
    \rowcolors{2}{mdalternatingtablerowcolor}{white}
    \begin{tabular}{!{\color{mdtableframecolor}\vrule}l!%
{\color{mdtableframecolor}\vrule}l!{\color{mdtableframecolor}\vrule}}
        \arrayrulecolor{mdtableframecolor}\hline
        \textbf{Distribution} & \textbf{Package name}\\
        \arrayrulecolor{mdtableframecolor}\hline
        Arch Linux & \mdinlinecode{ncurses}\\
        Debian & \mdinlinecode{ncurses-bin}\\
        Fedora & \mdinlinecode{ncurses}\\
        openSUSE & \mdinlinecode{ncurses-utils}\\
        Ubuntu & \mdinlinecode{ncurses-bin}\\
        \arrayrulecolor{mdtableframecolor}\hline
    \end{tabular}
\end{table}}
\end{mdcodeblock}

{\linespread{1.9}
\begin{table}[H]
    \rowcolors{2}{mdalternatingtablerowcolor}{white}
    \begin{tabular}{!{\color{mdtableframecolor}\vrule}l!%
{\color{mdtableframecolor}\vrule}l!{\color{mdtableframecolor}\vrule}}
        \arrayrulecolor{mdtableframecolor}\hline
        \textbf{Distribution} & \textbf{Package name}\\
        \arrayrulecolor{mdtableframecolor}\hline
        Arch Linux & \mdinlinecode{ncurses}\\
        Debian & \mdinlinecode{ncurses-bin}\\
        Fedora & \mdinlinecode{ncurses}\\
        openSUSE & \mdinlinecode{ncurses-utils}\\
        Ubuntu & \mdinlinecode{ncurses-bin}\\
        \arrayrulecolor{mdtableframecolor}\hline
    \end{tabular}
\end{table}}

Table created using only \mdinlinecode{tabu}:

\begin{mdcodeblock}
\taburowcolors[2] {white .. mdalternatingtablerowcolor}
\taburulecolor{mdtableframecolor}
\tabulinesep=^1.2ex
\tabulinesep=_0.75ex
\arrayrulewidth=0.4pt
\everyrow{\hline}
\begin{tabu}{|l|l|}
\rowfont[c]\bfseries
Distribution & Package name\\
Arch Linux & \mdinlinecode{ncurses}\\
Debian & \mdinlinecode{ncurses-bin}\\
Fedora & \mdinlinecode{ncurses}\\
openSUSE & \mdinlinecode{ncurses-utils}\\
Ubuntu & \mdinlinecode{ncurses-bin}\\
\end{tabu}
\end{mdcodeblock}

\taburowcolors[2] {white .. mdalternatingtablerowcolor}
\taburulecolor{mdtableframecolor}
\tabulinesep=^1.2ex
\tabulinesep=_0.75ex
\arrayrulewidth=0.4pt
\everyrow{\hline}
\begin{tabu}{|l|l|}
\rowfont[c]\bfseries
Distribution & Package name\\
Arch Linux & \mdinlinecode{ncurses}\\
Debian & \mdinlinecode{ncurses-bin}\\
Fedora & \mdinlinecode{ncurses}\\
openSUSE & \mdinlinecode{ncurses-utils}\\
Ubuntu & \mdinlinecode{ncurses-bin}\\
\end{tabu}

\newpage

Also created with \mdinlinecode{tabu}:

\begin{mdcodeblock}
\begin{table}[H]
\captionsetup{justification=justified, singlelinecheck=false}
\newcolumntype{C}[1]{>{\centering}m{#1}}
\taburowcolors[2] {white .. mdalternatingtablerowcolor}
\taburulecolor{mdtableframecolor}
\tabulinesep=^1.2ex
\tabulinesep=_0.75ex
\arrayrulewidth=0.4pt
\everyrow{\hline}
\begin{tabu}{|l|C{2.5cm}|C{2.5cm}|p{5.5cm}|}
\rowfont[c]\bfseries
Day & Minimum Temperature & Maximum Temperature & Summary \\
Monday & 11 \textdegree{}C & 22 \textdegree{}C & A clear day with lots
of sunshine. However, the strong breeze will bring down the temperatures. \\
Tuesday & ~9 \textdegree{}C & 19 \textdegree{}C & Cloudy with rain,
across many northern regions. Clear spells across most of Scotland and
Northern Ireland, but rain reaching the far northwest. \\
Wednesday & 10 \textdegree{}C & 21 \textdegree{}C & Rain will still linger
for the morning. Conditions will improve by early afternoon
and continue throughout the evening. \\
\end{tabu}
\caption{Weather forecast for the next 3 days.}
\end{table}
\end{mdcodeblock}

\begin{table}[H]
\captionsetup{justification=justified, singlelinecheck=false}
\newcolumntype{C}[1]{>{\centering}m{#1}}
\taburowcolors[2] {white .. mdalternatingtablerowcolor}
\taburulecolor{mdtableframecolor}
\tabulinesep=^1.2ex
\tabulinesep=_0.75ex
\arrayrulewidth=0.4pt
\everyrow{\hline}
\begin{tabu}{|l|C{2.5cm}|C{2.5cm}|p{5.5cm}|}
\rowfont[c]\bfseries
Day & Minimum Temperature & Maximum Temperature & Summary \\
Monday & 11 \textdegree{}C & 22 \textdegree{}C & A clear day with lots
of sunshine. However, the strong breeze will bring down the temperatures. \\
Tuesday & ~9 \textdegree{}C & 19 \textdegree{}C & Cloudy with rain,
across many northern regions. Clear spells across most of Scotland and
Northern Ireland, but rain reaching the far northwest. \\
Wednesday & 10 \textdegree{}C & 21 \textdegree{}C & Rain will still linger
for the morning. Conditions will improve by early afternoon
and continue throughout the evening. \\
\end{tabu}
\caption{Weather forecast for the next 3 days.}
\end{table}

Text source: \url{http://en.wikibooks.org/wiki/LaTeX/Tables#Text_wrapping_in_tables}

\mdsimplehrule

\section{Blockquotes}

Blockquotes with the \mdinlinecode{mdblockquote} environment:

\begin{mdblockquote}
The Three Laws of Robotics (often shortened to The Three Laws or Three Laws) are a set of rules devised by the science fiction author Isaac Asimov.

The rules were introduced in his 1942 short story ``Runaround'', although they had been foreshadowed in a few earlier stories. The Three Laws are:

\begin{mdenumerate}
\item A robot may not injure a human being or, through inaction, allow a human being to come to harm.
\item A robot must obey the orders given to it by human beings, except where such orders would conflict with the First Law.
\item A robot must protect its own existence as long as such protection does not conflict with the First or Second Law.
\end{mdenumerate}
\end{mdblockquote}

Source: \href{http://en.wikipedia.org/wiki/Three_Laws_of_Robotics}{Three Laws of Robotics -- Wikipedia, the free encyclopedia}

\mdsimplehrule

\section{Horizontal rules}

Fancy horizontal rule with \mdinlinecode{\textbackslash{}mdfancyhrule}:

\mdfancyhrule

Simple horizontal rule with \mdinlinecode{\textbackslash{}mdsimplehrule}:

\mdsimplehrule

\section{etc.}

Table of contents with \mdinlinecode{\textbackslash{}mdtableofcontents}.

List of figures and tables are the same as usual (\mdinlinecode{\textbackslash{}listoffigures}, \mdinlinecode{\textbackslash{}listoftables}).

\mdsimplehrule

\end{document}